% 第 14 题：粒子在电容器与磁场中的运动装置（TikZ + circuitikz）
% 左：电源 E0（内阻 r0）与滑动变阻器 R（最大 2r0），滑片在中点；
%     电容两极板并联在变阻器下半部分（滑片与下端）两端；
% 中：水平平行板电容器（板间距 d、板长 √3d），粒子由 a 进入、经 b 出、由 c 回；
% 右：垂直纸面向里的匀强磁场区域（× 标记）。
\begin{tikzpicture}[font=\small, scale=0.8, every node/.style={transform shape=false}]
  % ---------- 左：电源、内阻与变阻器电路 ----------
  \draw (0,2.2) to[battery, l=$E_{0}$] (0,0);
  \draw (0,2.2) to[R, l=$r_{0}$] (-2.2,2.2);
  % 滑动变阻器（竖直，最大 2r0，滑片在中点 (-2.2,0.8)）
  \draw (-2.2,2.2) to[vR, l=$R$] (-2.2,-0.6);
  \fill (-2.2,0.8) circle (0.06);
  % 变阻器下端 → 电池负极
  \draw (-2.2,-0.6) -- (0,-0.6) -- (0,0);
  % 滑片（中点）→ 上极板；变阻器下端 → 下极板
  \draw (-2.2,0.8) -- (0.3,0.8) -- (0.3,3.4) -- (0.6,3.4);
  \draw (-2.2,-0.6) -- (0.3,-0.6) -- (0.3,1.0) -- (0.6,1.0);
  % ---------- 中：平行板电容器（上极板 y=3.4，下极板 y=1.0，板长 √3d≈4.2）----------
  \draw[very thick] (0.6,3.4) -- (4.8,3.4);
  \draw[very thick] (0.6,1.0) -- (4.8,1.0);
  \draw[<->] (5.1,1.0) -- (5.1,3.4);
  \node[right] at (5.1,2.2) {$d$};
  % 粒子轨迹：a（左侧中点）→ b（下板右缘）→ 磁场圆弧 → c（上板右缘）
  \draw[dashed, -{Stealth}] (0.6,2.2) .. controls (1.6,2.1) and (3.4,1.6) .. (4.8,1.0);
  \draw[dashed, -{Stealth}] (4.8,1.0) .. controls (5.9,0.7) and (6.4,1.4) .. (5.9,2.2)
        .. controls (5.55,2.85) and (5.1,3.1) .. (4.8,3.4);
  % 关键点标注
  \fill (0.6,2.2) circle (0.07);
  \fill (4.8,1.0) circle (0.07);
  \fill (4.8,3.4) circle (0.07);
  \node[left] at (0.6,2.2) {$a$};
  \node[below right] at (4.8,1.0) {$b$};
  \node[above right] at (4.8,3.4) {$c$};
  % 板长标注
  \draw[<->] (0.6,3.8) -- (4.8,3.8);
  \node[above] at (2.7,3.8) {$\sqrt{3}\,d$};
  % ---------- 右：匀强磁场区域（× 标记）----------
  \foreach \x in {6.8,7.6,8.4,9.2} {
    \foreach \y in {0.7,1.5,2.3,3.1} {
      \node[red!70!black] at (\x,\y) {$\times$};
    }
  }
  \draw[dashed, rounded corners] (6.4,0.3) rectangle (9.7,3.8);
  \node at (8.05,-0.35) {$B$（垂直纸面向里）};
\end{tikzpicture}
